\section{Game-Based Definitions}
As one can see from the minimal differences between the \emph{passive} security definition and the \emph{active} security definition for PKE from Section~\ref{sec:overview:pke}, definitions of security can be modular with respect to a pattern that is based on a small core.
We will use the example of \emph{Key Encapsulation Mechanisms}~(KEM) to introduce a methodology for defining security systematically.
This methodology is closely related to a concept called \emph{Indistinguishability up to Correctness} by Rogaway and Zhang~\cite{C:RogZha18}.

\subsection{Syntax}
The first step to formally consider a cryptographic primitive is defining its \emph{syntax}.
The syntax specifies the algorithms of a primitive by defining their inputs and outputs.

A key encapsulation mechanism $\KEM=(\KEMgen,\KEMenc,\KEMdec)$ is a tuple of three algorithms with encapsulation key space~$\eksp$, decapsulation key space~$\dksp$, ciphertext space~$\csp$, and symmetric key space~$\ksp$:
\begin{itemize}
    \item $\KEMgen: \emptyset \tor \dksp\times\eksp$\\ 
    Generates a key pair of decapsulation key~$\dk\in\dksp$ and encapsulation key~$\ek\in\eksp$
    \item $\KEMenc: \eksp \tor \ksp\times\csp$\\
    Encapsulates a message~$k\in\ksp$ with encapsulation key~$\ek\in\eksp$ in ciphertext~$c\in\csp$
    \item $\KEMdec: \dksp\times\csp \to \ksp$\\
    Decapsulates key~$k\in\ksp$ with decapsulation key~$\dk\in\dksp$ from ciphertext~$c\in\csp$
\end{itemize}

Intuitively a KEM can be considered a special case of PKE that is limited to always encrypt (aka.\ encapsulate) fresh symmetric keys.
Instead of taking the symmetric keys as input, the encapsulation algorithm of a KEM generates them internally.

\paragraph{Construction: ElGamal KEM}
For illustration, we give the ElGamal KEM as a simple example of a KEM in Figure~\ref{fig:kem:const:elgamal}.

\begin{figure}[!ht]
    \centering
    \nicoresetlinenr%
    \fbox{%
        \scalebox{\codescalefactor}{%
            \input{figures/kem_const_elgamal}%
        }%
    }
    \caption{%
        ElGamal key encapsulation mechanism~$\KEM$~\cite{ElGamal85}.
    }
    \label{fig:kem:const:elgamal}
\end{figure}

\subsection{Correctness}
Before specifying the expected security guarantees, we define the \emph{correctness} requirements.
Typically, these requirements capture the functionality that a primitive should offer when being executed in an \emph{honest} environment in which all traffic is delivered as intended.
The literature on distributed systems developed two complementary notions of correctness:
\begin{itemize}
    \item \textbf{Safety:} Nothing ``bad'' should happen (e.g., encapsulation and decapsulation for the same ciphertext should not compute different symmetric keys)
    \item \textbf{Liveness:} The ``good'' event always happens (e.g., encapsulation and decapsulation always output a non-trivial key)
\end{itemize}

A correctness definition in the sense of \emph{safety} typically suffices for the ultimate goal of defining and analyzing security.
Furthermore, specifying \emph{liveness} requirements for sophisticated primitives sometimes leads to complex, incomprehensible definitions.

For the simple case of a KEM, a correctness definition in the flavor of safety would be:
\[
\Pr[\KEMdec(\dk,c)\in\{k,\bot\}\mid (\dk,\ek)\getsr\KEMgen,(k,c)\getsr\KEMenc(\ek))\alert{=m}]=1
\]
\alert{[JR] What does ``$=m$'' mean? I removed it in the next two formulas.}
and a correctness definition in the flavor of liveness would be:
\[
\Pr[\KEMdec(\dk,c)\neq\bot\mid (\dk,\ek)\getsr\KEMgen,(k,c)\getsr\KEMenc(\ek))]=1\text{,}
\]
which, in combination, leads to the standard correctness definitions for a~$\KEM$:
\[
\Pr[\KEMdec(\dk,c)=k\neq\bot\mid (\dk,\ek)\getsr\KEMgen,(k,c)\getsr\KEMenc(\ek))]=1\text{.}
\]

For more sophisticated primitives and protocols, one can tame complexity of definitions by specifying a correctness game.
Following the safety idea, a key encapsulation mechanism $\KEM$ is correct if $\Pr[\CORR_{\KEM}(\advA)=0]=1$ for all adversaries~$\advA$, where game~$\CORR$ is defined in Figure~\ref{fig:kem:corr}.

\begin{figure}[!ht]
    \centering
    \nicoresetlinenr%
    \fbox{%
        \scalebox{\codescalefactor}{%
            \input{figures/kem_corr}%
        }%
    }
    \caption{%
        Game $\CORR$ for key encapsulation mechanism~$\KEM$.
    }
    \label{fig:kem:corr}
\end{figure}

\subsection{Adversarial Capabilities}
Defining \emph{security} can be divided into three steps:
modeling the \emph{capabilities of adversaries},
specifying the \emph{security goal}, and
identifying attack strategies that are \emph{trivially} successful against every possible construction.

The former two steps are typically influenced by human intuition regarding the purpose of the considered primitive and how this primitive is used in practice.
This influence of human intuition can lead to ambiguities.
The latter step, identifying trivial attack strategies, is almost entirely determined by all remaining definitional steps, which reduces ambiguity.

\subsubsection{Typical Capabilities}
We begin with considering typical adversarial capabilities against KEMs:
\begin{itemize}
    \item \textbf{Chosen-Plaintext Attacks:}
    The adversary can see the ciphertexts of encapsulated keys
    \item \textbf{Chosen-Ciphertext Attacks:}
    The adversary can see the decapsulated keys for chosen ciphertexts
    \item \textbf{Secret Key Corruption:}
    The adversary can learn the entire decapsulation key
    \item \textbf{Leakage of Information during Algorithm Execution:}
    The adversary can learn secret bits of variables processed by the encapsulation or decapsulation algorithm
    \item \textbf{Subversion of Algorithms:}
    The adversary can modify the encapsulation or decapsulation algorithm
    \item Etc.
\end{itemize}

\paragraph{Chosen-Plaintext Attacks}
The most standard adversarial capabilities that are considered realistic and defendable for KEMs are chosen-plaintext and chosen-ciphertext attacks (CPA resp.\ CCA).

Consider the adversarial capabilities shown in Figure~\ref{fig:adv_capabilities}.
A passive adversary can see the encapsulation key~$\ek$ and the ciphertext~$c$ (dashed boxes).
Observing the values~$\ek$ and~$c$ is motivated by simple eavesdropping attacks, which models CPA.

\begin{figure}[!ht]
    \centering
    %!TEX root=../main.tex
\begin{tabular}{lcl}
    \underline{Alice} & \underline{$\advA$} & \underline{Bob}\\
    $(k,c)\getsr\KEMenc(\ek)$ & \arrbox{2cm}{$\xleftarrow{\dbox{$\ek$}}$} & $(\dk,\ek)\getsr\KEMgen$\\
    & $\xrightarrow{\dbox{$c$}}$ & $k\gets\KEMdec(\dk,c)$\\
    & \textcolor{redb}{$\xrightarrow{c'}$} & $\dbox{$k'$}\gets\KEMdec(\dk,c')$
\end{tabular}%
    %Provide alternative text for the caption, so that we can use \dbox
    \caption[A sketch of common adversarial capabilities]{%
        A sketch of common adversarial capabilities for \dbox{passive} and \textcolor{redb}{active} adversaries.
    }
    \label{fig:adv_capabilities}
\end{figure}

\paragraph{Chosen-Ciphertext Attacks}
The fact that~$k'$ is visible is motivated by active adversaries who can send~$c'$ and observe certain \emph{side-channels} to learn information about~$k'$.
For example, the adversary can measure the time it takes to decapsulate~$c'$ or send an invalid~$c'$ and observe the error message.
To model this capability we allow the adversary to choose ciphertexts for decapsulation and provide it with the decapsulation oracle~(CCA).
This decapsulation oracle models all forms of leaked side-channel information that can be obtained by observing the decapsulation.

The game $\INDCCA$, defined in Figure~\ref{fig:kem:ind}, captures chosen-ciphertext attacks as oracle $\Odec$ is provided;
if querying this oracle is forbidden, only chosen-plaintext attacks are modeled.

\begin{figure}[!ht]
    \centering
    \nicoresetlinenr%
    \fbox{%
        \scalebox{\codescalefactor}{%
            \input{figures/kem_ind}%
        }%
    }
    \caption{%
        Games $\INDCCA$ for key encapsulation mechanism~$\KEM$.
    }
    \label{fig:kem:ind}
\end{figure}

\subsection{Security Goal}
Specifying security goals is our second step.
It describes what the adversary is trying to break and what the cryptographic primitive is trying to protect/prevent.
The security goal is typically denoted in the first part of the name of the security game (e.g. \emph{IND}-CPA or \emph{OW}-CCA).

\subsubsection{Typical Goals}
For KEMs, the typical security goals are either \emph{indistinguishability} or \emph{one-way security} of the encapsulated keys.

\paragraph{One-Way Security}
For a KEM, one-way~(OW) security means that it must be hard to learn the full encapsulated key~$k$.
Thus, an attacker is only successful if all bits of the key are derived correctly.
The security game $\OWCCA$ is defined in Figure~\ref{fig:kem:ow:cca} and the $\OWCPA$ game follows from it by excluding the decapsulation oracle.

\begin{figure}[!ht]
    \centering
    \nicoresetlinenr%
    \fbox{%
        \scalebox{\codescalefactor}{%
            \input{figures/kem_ow_cca.tex}%
        }%
    }
    \caption{%
        One-Way CCA security game for $\KEM$.
    }
    \label{fig:kem:ow:cca}
\end{figure}

\paragraph{Indistinguishability}
As we have seen in Section~\ref{sec:overview:pke}, we sometimes want \emph{indistinguishability} of the secret encrypted information from some other chosen or randomly sampled bit strings.
The reason for this is that it gives us a stronger security guarantee. 
Namely, if a construction is IND secure, the adversary is unable to learn a single bit of the secret.
The $\INDCCA$ game is defined in Figure~\ref{fig:kem:ind}

\subsection{Trivial Winning Strategies}\label{sec:kem:trivial_attacks}
So far, we have done some unexplained bookkeeping in our security games. The reason for that are \emph{trivial winning strategies}.
These are strategies or attacks that work against every correct construction, simply because correctness enforces that the constructions behave \emph{meaningfully}.
For our KEM, there are two trivial winning strategies we need to guard against (Figure~\ref{kem:triv}):
The decapsulation oracle is used to learn the key from the challenge ciphertext; this key is either directly used as the game's solution ($\OWCCA$), or it is compared with the challenge key to win the game ($\INDCCA$).

\begin{figure}[!ht]%
    \centering
    \input{figures/kem_triv.tex}
    \caption{Trivial winning strategies against $\OWCCA$ and $\INDCCA$ $\KEM$ games.}
    \label{kem:triv}
\end{figure}

Now that the trivial winning conditions are defined, we can specify the advantage terms for each game.
For the OW-CPA and OW-CCA games that is:
\begin{align*}
    \Adv_\KEM^\owcpa(\advA)\coloneqq&\Pr[\OWCPA_{\KEM}(\advA)=1]\\
    \Adv_\KEM^\owcca(\advA)\coloneqq&\Pr[\OWCCA_{\KEM}(\advA)=1]\text{.}
\end{align*}
The advantages for the IND games are defined as:
\begin{align*}
    \Adv_\KEM^\indcpa(\advA)\coloneqq&\left|\Pr[\INDCPA_{\KEM}^0(\advA)=1]-\Pr[\INDCPA_{\KEM}^1(\advA)=1]\right|\\
    \Adv_\KEM^\indcca(\advA)\coloneqq&\left|\Pr[\INDCCA_{\KEM}^0(\advA)=1]-\Pr[\INDCCA_{\KEM}^1(\advA)=1]\right|\text{.}
\end{align*}

To validate that all trivial winning conditions are caught and forbidden, it suffices to find a secure construction and prove that it satisfies the definition.

\paragraph{Multi-Instance Security}
The security definitions introduced so far only cover a single victim (or a pair of victims) that are attacked by an adversary.
In contrast, in the real-world an adversary can interact with many victims and break the security of only one of them.
To account for this, one can define a multi-instance security game (Figure~\ref{fig:kem:ind:mi:corrupt}).
This game is called \emph{multi-instance} because it captures a setting in which many instances are running concurrently.
This can be the case in TLS for example, where many connections are established at the same time.
If we were to use the IND-CCA game from Figure~\ref{fig:kem:ind}, then we would have to use a, so called, \emph{hybrid argument} to prove security.
The hybrid argument can be used to prove security for multiple instances, by showing step by step that each instance is secure and then showing the indistinguishability between each step.
Doing so, however, yields a factor~$q$ in the final security bound, where $q$ is the number of instances.
To avoid this degeneration of the security bound, we can use the multi-instance game, for which sometimes tighter security bounds exist.

\begin{figure}[!ht]
    \centering
    \nicoresetlinenr%
    \fbox{%
        \scalebox{\codescalefactor}{%
            \input{figures/kem_ind_mi_corrupt}%
        }%
    }
    \caption{%
        Multi-instance games $\INDCCA$ for key encapsulation mechanism~$\KEM$.
    }
    \label{fig:kem:ind:mi:corrupt}
\end{figure}